% GENERATED FILE. Edit canonical lessons, then run npm run generate:books.
% canonical-source-hash: fnv1a64:fc4387cd5852ab86
% canonical-lessons: IT-C321-presepe, IT-C321-panettone, IT-C321-cenone, IT-C321-addobbare, IT-C321-tombola

\chapter{Le feste d'inverno}
\label{ch:it-le-feste-d-inverno}
\hlchaptermodality{321}

\begin{chapteropening}
\textbf{By the end of this chapter:} \emph{I can describe the winter holidays: the nativity scene, the festive cake, the big dinner, decorating and the family game.}

Complete the last lesson of chapter 321: I can describe the winter holidays: the nativity scene, the festive cake, the big dinner, decorating and the family game.
\end{chapteropening}

\section[il presepe]{il presepe — the nativity scene, the crib}

\label{lesson:IT-C321-presepe}



\begin{quote}
Before the new one: say the Italian for the name day, then the Italian for the saint.
\end{quote}

\subsection*{You'll want to know: il presepe}
\textbf{il presepe} — \textquotedblleft{}the nativity scene, the crib\textquotedblright{}.

\begin{quote}
Ogni anno facciamo il presepe in salotto.
\end{quote}

\begin{quote}
\emph{Every year we set up the nativity scene in the living room.}
\end{quote}

Also \emph{il presepio}.

\begin{cousinweb}
\emph{presepe} is a learned borrowing from Latin \emph{praesaepe} \textquotedblleft{}a stall, a manger\textquotedblright{} (EtymDB).
\end{cousinweb}

\begin{grammarlens}[title={What you've built}]
One more word for telling, explaining and describing at B1.
\end{grammarlens}

\subsection*{Your turn}
\begin{itemize}
\raggedright
  \item \emph{Say it:} \emph{il presepe}
  \item \emph{Say it:} \emph{il presepe}, once more
  \item \emph{Say it:} a sentence of your own with \emph{il presepe}
  \item \emph{Recall:} say the Italian for the name day, then the Italian for the saint, then say \emph{il presepe} again
\end{itemize}

\begin{tcolorbox}[breakable,colback=teal!4,colframe=teal!35!black,title={Before you move on}]
What does il presepe mean? (\textquotedblleft{}The nativity scene, the crib\textquotedblright{}.) Say it once more.
\end{tcolorbox}
\section[il panettone]{il panettone — the panettone, the tall festive cake}

\label{lesson:IT-C321-panettone}



\begin{quote}
Before the new one: say the Italian for the saint, then the Italian for the nativity scene.
\end{quote}

\subsection*{You'll want to know: il panettone}
\textbf{il panettone} — \textquotedblleft{}the panettone, the tall festive cake\textquotedblright{}.

\begin{quote}
Il panettone ha l'uvetta e i canditi.
\end{quote}

\begin{quote}
\emph{Panettone has raisins and candied peel.}
\end{quote}

A sweet bread from the north, now eaten everywhere in winter.

\begin{cousinweb}
\emph{panettone} comes from Lombard dialect \emph{panaton} (EtymDB), built on \emph{pane} \textquotedblleft{}bread\textquotedblright{}, from Latin \emph{pānis}.
\end{cousinweb}

\begin{grammarlens}[title={What you've built}]
One more word for telling, explaining and describing at B1.
\end{grammarlens}

\subsection*{Your turn}
\begin{itemize}
\raggedright
  \item \emph{Say it:} \emph{il panettone}
  \item \emph{Say it:} \emph{il panettone}, once more
  \item \emph{Say it:} a sentence of your own with \emph{il panettone}
  \item \emph{Recall:} say the Italian for the saint, then the Italian for the nativity scene, then say \emph{il panettone} again
\end{itemize}

\begin{tcolorbox}[breakable,colback=teal!4,colframe=teal!35!black,title={Before you move on}]
What does il panettone mean? (\textquotedblleft{}The panettone, the tall festive cake\textquotedblright{}.) Say it once more.
\end{tcolorbox}
\section[il cenone]{il cenone — the big festive dinner}

\label{lesson:IT-C321-cenone}



\begin{quote}
Before the new one: say the Italian for the nativity scene, then the Italian for the panettone.
\end{quote}

\subsection*{You'll want to know: il cenone}
\textbf{il cenone} — \textquotedblleft{}the big festive dinner\textquotedblright{}.

\begin{quote}
Il cenone di fine anno è a casa nostra.
\end{quote}

\begin{quote}
\emph{The big New Year's Eve dinner is at our house.}
\end{quote}

The ending \emph{-one} makes things big: \emph{cena}, \emph{cenone}.

\begin{cousinweb}
\emph{cenone} is \emph{cena} \textquotedblleft{}dinner\textquotedblright{} with \emph{-one} \textquotedblleft{}big\textquotedblright{}; \emph{cena} is inherited from Latin \emph{cēna} (Lewis \& Short).
\end{cousinweb}

\begin{grammarlens}[title={What you've built}]
One more word for telling, explaining and describing at B1.
\end{grammarlens}

\subsection*{Your turn}
\begin{itemize}
\raggedright
  \item \emph{Say it:} \emph{il cenone}
  \item \emph{Say it:} \emph{il cenone}, once more
  \item \emph{Say it:} a sentence of your own with \emph{il cenone}
  \item \emph{Recall:} say the Italian for the nativity scene, then the Italian for the panettone, then say \emph{il cenone} again
\end{itemize}

\begin{tcolorbox}[breakable,colback=teal!4,colframe=teal!35!black,title={Before you move on}]
What does il cenone mean? (\textquotedblleft{}The big festive dinner\textquotedblright{}.) Say it once more.
\end{tcolorbox}
\section[addobbare]{addobbare — to decorate (a room, a tree)}

\label{lesson:IT-C321-addobbare}



\begin{quote}
Before the new one: say the Italian for the panettone, then the Italian for the big festive dinner.
\end{quote}

\subsection*{You'll want to know: addobbare}
\textbf{addobbare} — \textquotedblleft{}to decorate (a room, a tree)\textquotedblright{}.

\begin{quote}
Abbiamo addobbato la casa con le luci.
\end{quote}

\begin{quote}
\emph{We decorated the house with lights.}
\end{quote}

\emph{Gli addobbi} are the decorations.

\begin{cousinweb}
\emph{addobbare} was borrowed from Old French \emph{adober} \textquotedblleft{}to equip\textquotedblright{} (EtymDB), from a Frankish verb meaning \textquotedblleft{}to strike\textquotedblright{}, as in dubbing a knight.
\end{cousinweb}

\begin{grammarlens}[title={What you've built}]
One more word for telling, explaining and describing at B1.
\end{grammarlens}

\subsection*{Your turn}
\begin{itemize}
\raggedright
  \item \emph{Say it:} \emph{addobbare}
  \item \emph{Say it:} \emph{addobbare}, once more
  \item \emph{Say it:} a sentence of your own with \emph{addobbare}
  \item \emph{Recall:} say the Italian for the panettone, then the Italian for the big festive dinner, then say \emph{addobbare} again
\end{itemize}

\begin{tcolorbox}[breakable,colback=teal!4,colframe=teal!35!black,title={Before you move on}]
What does addobbare mean? (\textquotedblleft{}To decorate (a room, a tree)\textquotedblright{}.) Say it once more.
\end{tcolorbox}
\section[la tombola]{la tombola — the tombola, the family bingo game}

\label{lesson:IT-C321-tombola}



\begin{quote}
Before the new one: say the Italian for the big festive dinner, then the Italian for to decorate.
\end{quote}

\subsection*{You'll want to know: la tombola}
\textbf{la tombola} — \textquotedblleft{}the tombola, the family bingo game\textquotedblright{}.

\begin{quote}
Dopo cena giochiamo a tombola.
\end{quote}

\begin{quote}
\emph{After dinner we play tombola.}
\end{quote}

Numbers from 1 to 90; \emph{fare tombola} is to fill the whole card.

\begin{cousinweb}
\emph{tombola} is usually traced to \emph{tombolare} \textquotedblleft{}to tumble\textquotedblright{}, from the numbers tumbling out of the bag; \emph{tombolare} itself is usually explained as a sound word (EtymDB has no entry).
\end{cousinweb}

\begin{grammarlens}[title={What you've built}]
One more word for telling, explaining and describing at B1.
\end{grammarlens}

\subsection*{Your turn}
\begin{itemize}
\raggedright
  \item \emph{Say it:} \emph{la tombola}
  \item \emph{Say it:} \emph{la tombola}, once more
  \item \emph{Say it:} a sentence of your own with \emph{la tombola}
  \item \emph{Recall:} say the Italian for the big festive dinner, then the Italian for to decorate, then say \emph{la tombola} again
\end{itemize}

\begin{tcolorbox}[breakable,colback=teal!4,colframe=teal!35!black,title={Before you move on}]
What does la tombola mean? (\textquotedblleft{}The tombola, the family bingo game\textquotedblright{}.) Say it once more.
\end{tcolorbox}
